\documentclass[10pt,a4paper]{amsart}
\usepackage[margin=19mm]{geometry}
\usepackage{amsmath,amssymb,mathtools}
\usepackage[scr=rsfs,cal=euler]{mathalfa}
\usepackage{xfrac}
\usepackage[unicode,hidelinks,bookmarks=false]{hyperref}
\newcommand{\ie}{i.e.\ }
\newcommand{\cf}{cf.\ }
\newcommand{\resp}{respectively\ }
\newcommand{\Subsection}{\subsection}
\providecommand{\nobreakdash}{\nobreak\hbox{-}}
\setlength{\emergencystretch}{2em}
\multlinegap0pt
\usepackage{bm}
\usepackage{bbm}
\usepackage{dsfont}
\usepackage{xspace}
\usepackage{cancel}
\usepackage{mathtools}
\usepackage{cleveref}
\usepackage[shortlabels]{enumitem}
\usepackage{amsthm}
\makeatletter
\@ifundefined{theorem}{\newtheorem{theorem}{Theorem}}{}
\@ifundefined{Proposition}{\newtheorem{Proposition}[theorem]{Proposition}}{}
\@ifundefined{lemma}{\newtheorem{lemma}[theorem]{Lemma}}{}
\makeatother
\newcommand{\D}{\mathrm D}
\newcommand{\R}{\mathbb R}
\newcommand{\Sph}{\mathbb S}
\newcommand{\dd}{\,\mathrm d}
\newcommand{\so}{\mathfrak{so}}
\newcommand{\weakstar}{\mathrel{\stackrel{*}{\rightharpoonup}}}
\title[The distance between homotopy classes of Sobolev maps on sph]{The distance between homotopy classes of Sobolev maps on spheres}
\author{Rupert L. Frank, Paata Ivanisvili}
\date{Source: arXiv:2606.29382v1 (CC BY 4.0)}
\begin{document}
\maketitle

\noindent\emph{Source and licence.} The distance between homotopy classes of Sobolev maps on spheres. Version arXiv:2606.29382v1 (CC BY 4.0), distributed under the Creative Commons Attribution 4.0 International licence, as recorded in the arXiv licence metadata for this exact version and shown on the abstract page https://arxiv.org/abs/2606.29382v1. Full rights notice: licenses/arxiv-2606.29382-CC-BY-4.0.txt.\par\smallskip

\noindent\textit{Editorial scope.} Counted unit: the Proposition whose source label is \texttt{eq:F-properties} in the article (source0.tex lines 731--746, source label \texttt{eq:F-properties}), delivered together with the article's own complete original proof of it (source0.tex lines 748--827, one \texttt{proof} environment of 80 lines). No mathematics is written, repaired or re-derived: statement and proof are the article's own text line for line. Required context retained uncounted: lem:affine-spanning (lines 588--590); lem:one-d-averaging (lines 625--632); lem:sphere-averaging (lines 647--653); lem:profile (lines 686--692). Every definition, display and label of those lines is retained unchanged. Omitted from this package and not redistributed: 1--587, 591--624, 633--646, 654--685, 693--730, 747--747, 828--1090. Nothing inside a retained excerpt is omitted or altered: every retained body is delivered exactly as the article's lines are written, so no omission receipt of any kind accompanies this package. Citations: the retained excerpts carry no citation command at all, so no bibliography entry of the article is transcribed and no reference is redistributed. Reference adaptation: none was needed and none was made; every reference inside the delivered unit resolves inside it, so no unresolved reference or citation is produced. Printed slips kept literal: no printed word, symbol, hypothesis or number of the counted statement or of its proof was altered. Layout and class: the article's own class is \texttt{amsart} with packages including fontenc, lmodern, microtype, amsmath, amssymb, mathtools, enumitem, hyperref, cleveref; the wrapper is the lane's pinned \texttt{amsart} editorial wrapper at 10pt on A4, which restates the theorem-like environments the retained text needs and adds the article's own macro definitions that the retained text needs (\texttt{\textbackslash D}, \texttt{\textbackslash R}, \texttt{\textbackslash Sph}, \texttt{\textbackslash dd}, \texttt{\textbackslash so}, \texttt{\textbackslash weakstar}), with extra wrapper packages (bm, bbm, dsfont, xspace, cancel, mathtools, cleveref, enumitem). The delivered numbering is the unit's own. Assets: no retained span contains a figure environment, a graphics-inclusion command, a PDF-page inclusion, a table or a TikZ picture; no image file is copied into this package, the package declares no asset dependency and no image file is redistributed with it. Licence: version arXiv:2606.29382v1 of this article is distributed under the Creative Commons Attribution 4.0 International licence, as recorded in the arXiv licence metadata for this exact version and shown on the abstract page https://arxiv.org/abs/2606.29382v1. A full rights notice is delivered at licenses/arxiv-2606.29382-CC-BY-4.0.txt. Characters: every retained excerpt is pure ASCII, so the delivered engine, XeLaTeX with its default Latin Modern text font, reports no missing character.
\begin{lemma}[Uniform affine spanning]\label[lemma]{lem:affine-spanning}
For every $p\in\Sph^n$, the set $\{Qp:Q\in\mathcal Q_\theta\}$ affinely spans $\R^{n+1}$.
\end{lemma}

\begin{lemma}[One-dimensional periodic averaging]\label[lemma]{lem:one-d-averaging}
Let $a\in L^\infty(\R)$ be one-periodic and $\bar a:=\int_0^1a(s)\,\dd s$.  Then
\[
  a(jt)\weakstar\bar a
  \quad\text{in }L^\infty((-1,1))
\]
as $j\to\infty$ through positive integers.
\end{lemma}

\begin{lemma}[Periodic averaging on $\Sph^n$]\label[lemma]{lem:sphere-averaging}
Under the assumptions of \Cref{lem:one-d-averaging},
\begin{equation}\label{eq:sphere-average}
  a(jx_{n+1})\weakstar\bar a
  \quad\text{in }L^\infty(\Sph^n).
\end{equation}
\end{lemma}

\begin{lemma}[Smooth periodic rotation profile]\label[lemma]{lem:profile}
Fix $\rho>0$.  There are a nonzero sufficiently small angle $\theta$, a smooth one-periodic map $B:\R\to\so(n+1)$, and pairwise disjoint intervals $I_Q\subset(0,1)$ of positive length, indexed by $Q\in\mathcal Q_\theta$, such that
\begin{enumerate}[label=\textup{(\roman*)}]
\item $\exp(B(t))=Q$ for $t\in I_Q$;
\item $|\exp(B(t))p-p|\le\rho$ for every $t\in\R$ and $p\in\Sph^n$.
\end{enumerate}
\end{lemma}

\begin{Proposition}[Oscillatory pinning]\label[Proposition]{prop:pinning}
Let $f_0\in C^\infty(\Sph^n;\Sph^n)$ and $\rho>0$.  There is a sequence $F_j\in C^\infty(\Sph^n;\Sph^n)$ such that
\begin{equation}\label{eq:F-properties}
  \deg F_j=\deg f_0,
  \qquad\|F_j-f_0\|_{L^\infty}\le\rho,
\end{equation}
and with the following property: if $g_j\in W^{1,n}(\Sph^n;\Sph^n)$ and
\begin{equation}\label{eq:bounded-difference}
  \sup_j\int_{\Sph^n}|\D(F_j-g_j)|^n\,\dd\sigma<\infty,
\end{equation}
then
\begin{equation}\label{eq:pinning-convergence}
  F_j-g_j\to0
  \quad\text{strongly in }L^n(\Sph^n;\R^{n+1}).
\end{equation}
\end{Proposition}

\begin{proof}
Let $\theta$ and $B$ be as in \Cref{lem:profile}, and define
\begin{equation}\label{eq:Fj-def}
  F_j(x):=\exp(B(jx_{n+1}))f_0(x).
\end{equation}
The claimed $L^\infty$ estimate on $F_j-f_0$ follows from \Cref{lem:profile}(ii).  The maps
\[
  H_s(x):=\exp(sB(jx_{n+1}))f_0(x),
  \qquad0\le s\le1,
\]
form a smooth homotopy from $f_0$ to $F_j$, so the degrees of $F_j$ and $f_0$ agree.

Let $g_j\in W^{1,n}(\Sph^n;\Sph^n)$ satisfy \eqref{eq:bounded-difference} and put
\[
  h_j:=F_j-g_j.
\]
Because both maps are sphere-valued,
\begin{equation}\label{eq:h-bound}
  |h_j|\le2
\end{equation}
and
\begin{equation}\label{eq:sphere-identity}
  2F_j\cdot h_j=|h_j|^2
  \qquad\text{a.e.}
\end{equation}
The bounds \eqref{eq:bounded-difference} and \eqref{eq:h-bound} make $(h_j)$ bounded in $W^{1,n}(\Sph^n;\R^{n+1})$. 

We are going to show that $h_j\to 0$ strongly in $L^n(\Sph^n;\R^{n+1})$, and we are going to show this by proving that any subsequence has a further subsequence that converges strongly to zero. 

In order to simplify the presentation, we do not reflect the choice of a subsequence in the notation. Since $n\ge2$, Rellich's compactness theorem gives, after taking a subsequence,
\begin{equation}\label{eq:h-L2}
  h_j\to h\quad\text{strongly in }L^2(\Sph^n)
\end{equation}
and we need to show that $h=0$.

Fix $Q\in\mathcal Q_\theta$ and let $\chi_Q$ be the one-periodic extension of the characteristic function of $I_Q$.  Where $\chi_Q(jx_{n+1})=1$, one has $F_j=Qf_0$, and therefore
\begin{equation}\label{eq:plateau-identity}
  \chi_Q(jx_{n+1})
  \bigl(2Qf_0\cdot h_j-|h_j|^2\bigr)=0.
\end{equation}
By \Cref{lem:sphere-averaging},
\begin{equation}\label{eq:chi-average}
  \chi_Q(jx_{n+1})\weakstar |I_Q|
  \quad\text{in }L^\infty(\Sph^n).
\end{equation}
Also, \eqref{eq:h-L2} implies
\begin{equation}\label{eq:nonlinear-L1}
  2Qf_0\cdot h_j-|h_j|^2
  \to2Qf_0\cdot h-|h|^2
  \quad\text{strongly in }L^1.
\end{equation}
Indeed,
\[
  \||h_j|^2-|h|^2\|_{L^1}
  \le\|h_j-h\|_{L^2}\|h_j+h\|_{L^2}\to0.
\]

Testing \eqref{eq:plateau-identity} against an arbitrary $\varphi\in L^\infty(\Sph^n)$ and using \eqref{eq:chi-average} and \eqref{eq:nonlinear-L1}, we obtain
\[
  |I_Q|\int_{\Sph^n}\varphi\,\bigl(2Qf_0\cdot h-|h|^2\bigr)\,\dd\sigma=0.
\]
Since $|I_Q|>0$, this implies
\begin{equation}\label{eq:limit-affine-equations}
  2Qf_0\cdot h=|h|^2
  \qquad\text{a.e.~on}\ \Sph^n \,.
\end{equation}
This holds for every $Q\in\mathcal Q_\theta$.

Since $\mathcal Q_\theta$ is finite, there is a set of full measure on which \eqref{eq:limit-affine-equations} holds for all $Q\in\mathcal Q_\theta$.  At such a point $x$, if $h(x)\ne0$, then every vector $Qf_0(x)$ lies in the proper affine hyperplane
\[
  \left\{y\in\R^{n+1}:y\cdot h(x)=\frac{|h(x)|^2}{2}\right\}.
\]
This contradicts \Cref{lem:affine-spanning}.  Hence $h=0$ a.e.

We have shown that every subsequence of $(h_j)$ has a further subsequence converging to zero strongly in $L^2$.  Therefore the full sequence converges to zero in $L^2$.  Since $|h_j|\le2$,
\[
  \|h_j\|_{L^n}^n\le 2^{n-2}\|h_j\|_{L^2}^2\to0,
\]
which proves \eqref{eq:pinning-convergence}.
\end{proof}

\end{document}
